\documentclass[]{article}

\usepackage{amsmath}
\usepackage{multirow}
\usepackage{natbib}
\usepackage{graphicx} 


\begin{document}

\subsection{Eulerian spectral roughness from multifractal and deterministic models}

The first step in our investigation was to characterize the Eulerian velocity fields by measuring their spectral roughness, $\xi$. This quantity, derived from the second-order structure function exponent $\zeta_2$ via the relation $\xi = 1 + \zeta_2$, provides the theoretical prediction for the Lagrangian transport exponent, $\alpha = 2/\xi$. We computed the structure function scaling exponents, $\zeta_p$, for the multifractal energy cascade models and the deterministic Lorenz-96 system.

The results are summarized in Figure~\ref{fig:image1}. For the non-intermittent, Kolmogorov-like cascade ($\mu=0.0$), the scaling exponents follow the theoretical linear relation $\zeta_p = p/3$. This yields a second-order exponent of $\zeta_2 = 2/3$ and a corresponding spectral roughness of $\xi = 5/3 \approx 1.667$, in exact agreement with classical Kolmogorov theory. As intermittency is introduced, the $\zeta_p$ curve becomes nonlinear and concave, consistent with multifractal theory. For mild intermittency ($\mu=0.15$), we measure $\xi = 1.683$, and for realistic intermittency ($\mu=0.28$), the roughness increases further to $\xi = 1.697$. In stark contrast, the chaotic Lorenz-96 system exhibits a much flatter scaling profile, yielding $\zeta_2 = 0.029$ and a spectral roughness of $\xi = 1.029$. This indicates a significantly smoother effective velocity field compared to the turbulent cascades.

These Eulerian measurements establish the baseline for our predictions. The theory posits that the increased roughness associated with intermittency should manifest as a decrease in the Lagrangian Lévy index $\alpha$, implying more pronounced anomalous transport with heavier-tailed displacement distributions.

\begin{figure}[h!]
    \centering
    \includegraphics[width=0.5\textwidth]{../input_files/plots/step_2_zeta_p_vs_p_2_1778360956.png}
    \caption{Scaling exponents $\zeta_p$ of the Eulerian velocity structure functions versus the moment order $p$ for several turbulence models. The empirical data for the pure Kolmogorov (K41) cascade align with the theoretical linear scaling $\zeta_p = p/3$. The introduction of intermittency via log-normal models (logN\_mild, logN\_realistic) results in the characteristic nonlinear, concave curvature, with the deviation from linearity increasing with the level of intermittency. In contrast, the deterministic Lorenz-96 model exhibits a significantly flatter scaling profile. These results demonstrate that intermittency systematically increases the second-order exponent $\zeta_2$, and therefore the Eulerian spectral roughness $\xi$, which is used to predict the nature of the Lagrangian transport.}
    \label{fig:image1}
\end{figure}

\subsection{Lagrangian temporal correlations}

A key assumption underlying the theoretical link between Eulerian roughness and fractional diffusion is that the Lagrangian velocity statistics are approximately white-in-time. This ensures that any observed anomalous transport arises from the long-range spatial correlations of the velocity field, rather than from temporal memory effects. To verify this condition, we computed the Lagrangian velocity autocorrelation function, $R_v(\tau)$, for the models used to generate tracer trajectories.

As shown in Figure~\ref{fig:image2}, the velocity fields in both the Kraichnan and Lorenz-96 models decorrelate rapidly. For the Kraichnan model, which is constructed to be delta-correlated in time, the autocorrelation function drops to zero within a single time step, as expected. The deterministic Lorenz-96 system also exhibits a very short correlation time ($\tau_c = 0.23$) relative to the total simulation duration ($T=500$). The satisfaction of this white-noise condition is critical, as it allows us to isolate the impact of the spatial structure of the flow on the Lagrangian dynamics.

\begin{figure}[h!]
    \centering
    \includegraphics[width=0.5\textwidth]{../input_files/plots/step_3_velocity_autocorr_3_1778361095.png}
    \caption{Lagrangian velocity autocorrelation functions, $R_v(\tau)$, for the Kraichnan model across various spectral roughness exponents ($\xi$) and for the deterministic Lorenz-96 model. The plot demonstrates that the velocity fields in both systems are effectively delta-correlated in time. For the Kraichnan model, the correlation drops instantaneously to zero, reflecting its construction. The Lorenz-96 model also exhibits a rapid decay to zero. This confirmation of the white-in-time approximation is a critical prerequisite, ensuring that observed anomalous diffusion is a consequence of the spatial velocity structure rather than temporal memory effects.}
    \label{fig:image2}
\end{figure}

\subsection{Pre-asymptotic constraints on Lagrangian transport in the Kraichnan model}

Having established the Eulerian predictions and verified the necessary temporal conditions, we next examined the Lagrangian transport statistics in the 3D Kraichnan model. This synthetic model serves as an ideal testbed, as the spectral roughness $\xi$ can be precisely controlled. The theory predicts superdiffusive transport, with the Mean Squared Displacement (MSD) scaling as $\langle \Delta x^2(\tau) \rangle \propto \tau^{2H}$ where the Hurst exponent $H = \xi/2$, and a displacement distribution characterized by a Lévy index $\alpha = 2/\xi$.

However, our empirical results show a significant departure from these asymptotic predictions. Figure~\ref{fig:image3} displays the MSD for several values of $\xi$. In all cases, the measured Hurst exponent remains close to $H \approx 0.5$, the value characteristic of normal, classical diffusion. This stands in stark contrast to the theoretical predictions of strong superdiffusion (e.g., $H=0.75$ for $\xi=1.5$).

\begin{figure}[h!]
    \centering
    \includegraphics[width=0.5\textwidth]{../input_files/plots/step_4_kraichnan_msd_4_1778361306.png}
    \caption{Log-log plot of the Mean Squared Displacement (MSD) versus time for Lagrangian tracers in the 3D Kraichnan model, shown for several values of the Eulerian spectral roughness exponent $\xi$. The legend displays the empirically measured Hurst exponent, $H$, for each case. Contrary to the theoretical prediction of anomalous diffusion where $H = \xi/2$, the measured scaling exponents consistently remain near $H \approx 0.5$, which is characteristic of normal diffusion. This reveals that the system is trapped in a pre-asymptotic state, where the finite size of the simulation's spectral domain prevents the tracers from experiencing the long-range correlations required to manifest the true anomalous scaling within the observed time frame.}
    \label{fig:image3}
\end{figure}

This near-diffusive behavior is further reflected in the shape of the Probability Density Functions (PDFs) of particle displacements, shown in Figure~\ref{fig:image5}. The PDFs are dominated by a Gaussian-like central bulk and lack the pronounced heavy tails that are the hallmark of Lévy-stable processes. To quantify this, we directly estimated the Lévy index $\alpha$ from the displacement data. As shown in the Hill plots in Figure~\ref{fig:image4}, the estimators for $\alpha$ are highly unstable and fail to converge to the theoretical prediction of $\alpha = 2/\xi$. Instead, they are strongly biased towards the Gaussian limit of $\alpha=2$.

Taken together, these results indicate that the system is trapped in a pre-asymptotic state. The finite spectral resolution of the simulation ($N_k=128$) and the finite simulation time prevent the tracer particles from sufficiently sampling the long-range spatial correlations of the velocity field. Consequently, the dynamics are dominated by the central part of the displacement distribution, and the theoretically predicted fractional behavior does not manifest.

\begin{figure}[h!]
    \centering
    \includegraphics[width=0.5\textwidth]{../input_files/plots/step_4_kraichnan_pdf_4_1778361306.png}
    \caption{Probability Density Functions (PDFs) of Lagrangian tracer displacements at t=40.0 in the Kraichnan model for different Eulerian spectral roughness exponents $\xi$. The distributions show a dominant central bulk, indicating that the system remains in a pre-asymptotic, near-Gaussian regime. This visualizes the impact of finite simulation size, which prevents the emergence of the heavy tails characteristic of the theoretically predicted asymptotic Lévy-stable statistics.}
    \label{fig:image5}
\end{figure}

\begin{figure}[h!]
    \centering
    \includegraphics[width=0.5\textwidth]{../input_files/plots/step_4_hill_stability_4_1778361306.png}
    \caption{Hill plots for the estimation of the Lévy index $\alpha$ from tracer displacements in the Kraichnan model across five values of the Eulerian spectral roughness, $\xi$. Each panel shows the estimated $\alpha$ versus the number of order statistics $k$ at two times, compared with the theoretical prediction $\alpha = 2/\xi$ (dashed line). For the anomalous superdiffusive cases ($\xi=1.5$ and $\xi=1.8$), the estimates fail to converge to the theory and are biased towards the Gaussian limit ($\alpha=2$). This demonstrates the pre-asymptotic trapping of the system, where finite-size effects cause the Gaussian bulk of the displacement distribution to dominate the expected fractional tails.}
    \label{fig:image4}
\end{figure}

\subsection{Renormalization group flow analysis}

To directly visualize the system's failure to reach the asymptotic fractional regime, we tracked the Renormalization Group (RG) flow of the effective Lévy exponent, $\alpha(\tau)$, as a function of the time lag $\tau$. This approach allows us to observe whether the system evolves from its short-time behavior towards the predicted fixed point. We extracted $\alpha(\tau)$ by fitting the empirical characteristic function of the displacements at each time lag.

The scaling of the characteristic function, $\phi(k, \tau)$, is shown in Figure~\ref{fig:image6}. For a Lévy-stable process, the negative logarithm of $\phi(k, \tau)$ should scale as $|k|^{\alpha(\tau)}$. The plot reveals that the slopes of the curves are nearly identical and parallel for all time lags, indicating that the effective exponent $\alpha(\tau)$ remains constant.

\begin{figure}[h!]
    \centering
    \includegraphics[width=0.5\textwidth]{../input_files/plots/step_5_char_func_scaling_5_1778361494.png}
    \caption{Scaling of the characteristic function $\phi(k, \tau)$ for different time lags $\tau$ in the Kraichnan model, plotted as the negative logarithm of the characteristic function versus wavenumber $|k|$ on a log-log scale. The slope of each line corresponds to the effective Lévy index $\alpha(\tau)$. The consistent, parallel slopes across all time lags ($\tau \in [0.01, 50.0]$) provide direct evidence that the system remains in a pre-asymptotic state with a near-constant exponent $\alpha(\tau) \approx 2.0$. This demonstrates that the Renormalization Group flow does not reach the predicted anomalous diffusive fixed point within the observed simulation time, trapping the dynamics in a near-Gaussian regime.}
    \label{fig:image6}
\end{figure}

The explicit RG flow is plotted in Figure~\ref{fig:image7}. This result provides the most direct evidence of pre-asymptotic trapping. The effective exponent $\alpha(\tau)$ remains robustly pinned near the Gaussian value of $\alpha=2.0$ across the entire range of observed time scales for both the $\xi=1.5$ and $\xi=1.8$ cases. The system shows no tendency to flow towards the theoretically predicted fixed points (dashed lines). This demonstrates that the crossover time required to transition into the anomalous regime is far greater than the total simulation time. The emergence of the fractional dynamics predicted by the Eulerian roughness is thus suppressed by the finite scale separation available in the simulation.

\begin{figure}[h!]
    \centering
    \includegraphics[width=0.5\textwidth]{../input_files/plots/step_5_rg_flow_alpha_5_1778361494.png}
    \caption{Renormalization Group (RG) flow of the effective Lévy index $\alpha(\tau)$ as a function of time lag $\tau$ for Kraichnan model simulations with spectral roughness $\xi=1.5$ and $\xi=1.8$. The empirical exponent, derived from a sliding-window analysis of the characteristic functions (black circles), remains pinned at the Gaussian value $\alpha \approx 2.0$ across the entire range of time lags. This demonstrates a failure to flow towards the theoretically predicted asymptotic fixed points $\alpha = 2/\xi$ (dashed lines), indicating that the simulation remains in a pre-asymptotic state where the crossover time to the fractional Lévy regime is much larger than the observation time.}
    \label{fig:image7}
\end{figure}

\subsection{The role of dimensionality: Topological trapping in 1D turbulence}

Finally, we investigated the validity of the $\alpha = 2/\xi$ mapping in a different physical context by analyzing tracer data from a one-dimensional synthetic turbulent flow. While the RG framework does not explicitly depend on dimensionality, the physical realization of transport can be strongly affected by topological constraints.

The results, shown in Figure~\ref{fig:image0}, reveal a stark breakdown of the theoretical mapping in 1D. For both the pure Kolmogorov and intermittent multifractal cases, the transport is strongly subdiffusive, with measured Hurst exponents of $H \approx 0.49$ and $H \approx 0.14$, respectively. The corresponding Lévy indices are $\alpha \approx 0.385$ and $\alpha \approx 0.439$. These values are drastically different from the superdiffusive prediction of $\alpha \approx 1.2$ derived from the Eulerian roughness.

This discrepancy arises from topological trapping. In a one-dimensional flow, tracer particles lack the freedom to move around regions of low velocity or stagnation points. They become trapped for extended periods, which introduces strong negative correlations in their displacements and suppresses the overall variance. This trapping mechanism dominates the transport dynamics, inducing a non-universal, subdiffusive behavior that is not captured by the theoretical framework relating transport to the field's spectral roughness. This finding highlights that the validity of the $\alpha = 2/\xi$ relation is contingent not only on sufficient scale separation but also on a spatial dimensionality that allows particles to freely sample the velocity field's fluctuations.

\begin{figure}[h!]
    \centering
    \includegraphics[width=0.5\textwidth]{../input_files/plots/step_6_kolmogorov_intermittency_6_1778360342.png}
    \caption{Analysis of 1D synthetic Kolmogorov tracer data for pure K41 (blue) and multifractal (red) turbulence, showing final displacement Probability Density Functions (left) and a comparison of the Lévy index $\alpha$ (right). The empirical results ($\alpha \approx 0.385$ for pure K41 and $\alpha \approx 0.439$ for multifractal) indicate severe subdiffusion, deviating significantly from the theoretical predictions for Lévy flights derived from the 3D Renormalization Group mapping ($\alpha \approx 1.2$). This discrepancy demonstrates that the theoretical mapping breaks down in 1D due to topological trapping effects that suppress particle transport.}
    \label{fig:image0}
\end{figure}

\end{document}
                